\documentclass[11pt]{article}

% --------------------
% Page setup
% --------------------
\usepackage[margin=1in]{geometry}
\usepackage{setspace}
\usepackage{parskip}

% --------------------
% Math packages
% --------------------
\usepackage{amsmath}
\usepackage{amssymb}
\usepackage{amsthm}
\usepackage{mathtools}
\usepackage{mathrsfs}
\usepackage{bm}

% --------------------
% Lists, tables, and formatting
% --------------------
\usepackage{enumitem}
\usepackage{booktabs}
\usepackage{xcolor}
\usepackage{framed}

% --------------------
% Diagrams and figures
% --------------------
\usepackage{tikz}
\usepackage{tikz-cd}
\usepackage{graphicx}
\usepackage{float}

% --------------------
% Links and references
% --------------------
\usepackage[hidelinks]{hyperref}
\usepackage[nameinlink,capitalise]{cleveref}

% --------------------
% Theorem environments
% --------------------
\theoremstyle{definition}
\newtheorem{definition}{Definition}[section]
\newtheorem{example}[definition]{Example}
\newtheorem{exercise}[definition]{Exercise}

\theoremstyle{plain}
\newtheorem{theorem}[definition]{Theorem}
\newtheorem{lemma}[definition]{Lemma}
\newtheorem{proposition}[definition]{Proposition}
\newtheorem{corollary}[definition]{Corollary}

\theoremstyle{remark}
\newtheorem{remark}[definition]{Remark}
\newtheorem{note}[definition]{Note}

% --------------------
% Common math shortcuts
% --------------------
\newcommand{\N}{\mathbb{N}}
\newcommand{\Z}{\mathbb{Z}}
\newcommand{\Q}{\mathbb{Q}}
\newcommand{\R}{\mathbb{R}}
\newcommand{\C}{\mathbb{C}}
\newcommand{\F}{\mathbb{F}}

\newcommand{\eps}{\varepsilon}
\newcommand{\set}[1]{\left\{#1\right\}}
\newcommand{\abs}[1]{\left|#1\right|}
\newcommand{\norm}[1]{\left\lVert#1\right\rVert}
\newcommand{\inner}[2]{\left\langle #1, #2 \right\rangle}
\newcommand{\paren}[1]{\left(#1\right)}
\newcommand{\bracket}[1]{\left[#1\right]}

\DeclareMathOperator{\im}{im}
\DeclareMathOperator{\id}{id}
\DeclareMathOperator{\Hom}{Hom}
\DeclareMathOperator{\Aut}{Aut}
\DeclareMathOperator{\Ker}{Ker}
\DeclareMathOperator{\rank}{rank}
\DeclareMathOperator{\Span}{span}

% --------------------
% Document info
% --------------------
\title{Math Study Guide}
\author{}
\date{\today}

\begin{document}

\maketitle
\tableofcontents
\newpage

\section{Topic Name}

\subsection{Definitions}

\begin{definition}


Orbit: 

How to find it:

\end{definition}

\begin{definition}

Stabilizer:

How to find it:

\end{definition}

\begin{definition}

Orbit-Stabilizer Theorem:

\end{definition}


\begin{definition}



\end{definition}

\begin{definition}

Lagrange's Theorem:

The size of the group is equal to the number of cosets times the size of each coset.

$$
\mid G \mid = \mid G/H \mid \cdot \mid H \mid
$$

The size of the quotient group of a group G by a subgroup H is equal to the number of cosets of H in G.

\end{definition}

\begin{definition}



\end{definition}





\subsection{Key Theorems}

\begin{theorem}

Chinese Remainder Theorem:

For coprime integers, you have the isomorphism:

$$Z_{n_1} \times Z_{n_2} \times \dots \times Z_{n_k} \cong Z_{n_1 n_2 \dots n_k}$$

\end{theorem}




\subsection{Examples}

\begin{example}

\end{example}

\subsection{Common Proof Patterns}

\begin{lemma}

\end{lemma}

\begin{proof}

\end{proof}

\subsection{Exercises}

\begin{exercise}

\end{exercise}

\end{document}
